\section{Conclusion}

We began with a simple observation: practitioners universally default to LoRA rank $r=16$ without principled justification. We end with an empirical scaling law: \textbf{optimal LoRA rank scales sub-linearly with model size}, following $r_{\text{opt}} \propto N^\alpha$ where $\alpha$ ranges from 0.30 to 0.82 depending on task type.

Larger models exhibit a phase transition in rank sensitivity---the penalty for suboptimal rank increases super-linearly with scale. Surprisingly, attention entropy \emph{decreases} with model size, suggesting larger models develop focused, specialized representations.

Moving forward, two directions are promising:
\begin{enumerate}
    \item \textbf{Task-calibrated $\alpha$}: Develop complexity metrics predicting the scaling exponent
    \item \textbf{Attention-focus scaling}: Test whether optimal rank correlates with inverse attention entropy
\end{enumerate}

The era of ad-hoc LoRA rank selection can end. With scaling laws as guides, efficient adaptation becomes predictable.
